\usepackage[section]{algorithm} % float wrapper for algorithms
\usepackage{algpseudocode}
\usepackage{amsfonts}
\usepackage{amsmath}
\usepackage{amssymb}
\usepackage{amsthm}
\usepackage{array}
\usepackage{cite} 
\usepackage{framed}
\usepackage{float} % table centering 
\usepackage{geometry}
\usepackage{graphicx}
\usepackage{hyperref} % creates hypertext links in pdf files (natbib compatible)
\usepackage{listings}
\usepackage{mathrsfs}
\usepackage{mathtools}
\usepackage{multirow}
\usepackage[numbers]{natbib} % round braces, sort multiple citations
\usepackage{parskip}
\usepackage{setspace}
\usepackage{subcaption}
\usepackage{xtab}


% The are a lot of things you might find useful and want to uncomment. 
% Most things you will want to uncomment or change will be near the top.

\newcommand{\keyword}[1]{\textit{{#1}}} % for keywords

% for standard maths operators
\newcommand{\abs}[1]{\ensuremath{\left| #1 \right|}}
\newcommand{\norm}[1]{\left \lVert #1 \right \rVert}
\newcommand{\Wtwo}{W_2}
\newcommand{\lvbracket}{[\mkern-4mu|}
\newcommand{\rvbracket}{{|\mkern-4mu]}}
\newcommand{\setbuild}{\, \left. \right| \,}

% standard maths symbols (you probably want to uncomment all of these)
\newcommand{\Z}{\ensuremath{\mathbb{Z}}}% integers
\newcommand{\N}{\ensuremath{\mathbb{N}}}% natural numbers
\newcommand{\R}{\ensuremath{\mathbb{R}}}% real numbers
\newcommand{\C}{\ensuremath{\mathbb{C}}}% complex numbers
\newcommand{\Q}{\ensuremath{\mathbb{Q}}} % rationals
\newcommand{\Psym}{\ensuremath{\mathbb{P}}} % Probability 
\newcommand{\Esym}{\ensuremath{\mathbb{E}}} % Expectation 

% Spaces
\newcommand{\ZeroN}[2]{\lvbracket #1, #2 \rvbracket}
\newcommand{\ZeroT}{[0, T]}
\newcommand{\ctsFunc}{C[0,T]}
\newcommand{\cdFunc}{D[0,T]}
\newcommand{\Ptwo}{\ensuremath{\mathcal{P}_2(\mathbb{R})}} 

%index 
\newcommand{\tinT}{\ensuremath{0 \leq t \leq T}}
\newcommand{\uind}[1]{\ensuremath{^{(#1)}}}

%For Probabiltiy and Expectation Operators 
\newcommand{\Prob}[1]{\ensuremath{\mathbb{P}}\left( #1 \right)} % probability 
\newcommand{\cProb}[2]{\mathbb{P}( #1 \, | \, #2)} % conditional probability
\newcommand{\Ex}[1]{\ensuremath{\mathbb{E}} \left[ #1 \right]} % expectation

% Processes 
\newcommand{\Xt}{\ensuremath{X_t}}
\newcommand{\Xs}{\ensuremath{X_s}}
\newcommand{\Xtdiscrete}{\ensuremath{X_i}}
\newcommand{\Wt}{\ensuremath{W_t}}
\newcommand{\Ws}{\ensuremath{W_s}}
\newcommand{\controlt}{\ensuremath{\mathfrak{a}_t}}
\newcommand{\controls}{\ensuremath{\mathfrak{a}_s}}
\newcommand{\Lawzero}{\ensuremath{\mathbb{P}_{X_0}}}
\newcommand{\Lawt}{\ensuremath{\mathbb{P}_{X_t}}}
\newcommand{\Laws}{\ensuremath{\mathbb{P}_{X_s}}}
\newcommand{\spro}[2]{#1_{#2}}
\newcommand{\Process}[2]{#1_{#2}}

% Empirical measure 
\newcommand{\empmeasure}{\ensuremath{\bar \mu_{\Xt}\uind{N}}}
\newcommand{\empmeasurediscrete}{\ensuremath{\bar \mu_{\Xtdiscrete}\uind{N}}}

% Functions 
\newcommand{\ulift}{\ensuremath{\tilde{u}}}
\newcommand{\minapprox}[1]{\varphi_\epsilon(#1)}
\newcommand{\decoupling}{\ensuremath{\partial_x V(t, x, \mu) + \int_\R \partial_\mu V(t, x', \mu)(v) \, \mu(dx')}}
\newcommand{\feedback}{\ensuremath {\mathfrak a}}
\newcommand{\nnfeedback}{\ensuremath{\mathcal{N}_\theta(t, x, \mu)}}

% Approximations 
\newcommand{\deltat}{\Delta t}

% Distributions 
\newcommand{\normal}[2]{\ensuremath{\mathcal{N}(#1, #2)}}
\newcommand{\uniform}[1]{\ensuremath{\text{Uniform}(0, #1)}}



% round braces, sort multiple citations
 % Citing the actual authors 
%\usepackage{hyperref}
%\usepackage{fancyhdr}
%\usepackage{mathrsfs}
%\usepackage{textcomp}
%\usepackage{color}
%\usepackage{graphicx}
%\usepackage{framed}
%\usepackage[vlined,boxed,commentsnumbered,algochapter]{algorithm2e}
%\usepackage{ifdraft}% perform operations conditional on the draft option

% \usepackage{mathptmx}
% \usepackage{mathpazo}
%\usepackage{amscd}
% \usepackage{xy}
% \usepackage{diagxy}
%\usepackage{diagrams}
%\usepackage{marginnote}
%\usepackage{rotating}
%\usepackage{multirow}
% \usepackage{polski}
% \usepackage[T1]{fontenc}

\usepackage{tikz, pgfplots}
\usetikzlibrary{tikzmark,fit}
\usetikzlibrary{matrix,shapes.arrows}
\usetikzlibrary{positioning}

% \usepackage[inline]{showlabels}
%\usepackage{booktabs}
% Date format


% New commands, operators and symbols

% Operators
%\DeclareMathOperator{\Sym}{Sym}% symmetric group
%\DeclareMathOperator{\Alt}{Alt}% alternating group
%\DeclareMathOperator{\Id}{Id}
%\DeclareMathOperator{\Hom}{Hom}
%\DeclareMathOperator{\Grp}{Grp}
%\DeclareMathOperator{\supp}{supp}
%\DeclareMathOperator{\fix}{fix}
%\DeclareMathOperator{\dep}{dep}
%\DeclareMathOperator{\lcm}{lcm}
%\DeclareMathOperator{\Aut}{Aut}
%\DeclareMathOperator{\Inn}{Inn}
%\DeclareMathOperator{\Out}{Out}
% \DeclareMathOperator{\dim}{dim}
%\DeclareMathOperator{\Syl}{Syl}
%\DeclareMathOperator{\Hall}{Hall}
%\DeclareMathOperator{\pCore}{pCore}
%\DeclareMathOperator{\Char}{char}
%\DeclareMathOperator{\Image}{Im}
%\DeclareMathOperator{\Ker}{Ker}
%\DeclareMathOperator{\Hcf}{hcf}
%\DeclareMathOperator{\GL}{GL}
%\DeclareMathOperator{\Pc}{Pc}
%\DeclareMathOperator{\Stab}{Stab}
%\DeclareMathOperator{\Orbit}{Orbit}

% Hyphenation fixes
%\newcommand{\letdash}[1]{$#1$\nobreakdash-\hspace{0pt}}% for n-element, k-transitive etc
%\newcommand{\numdash}{\nobreakdash--}

% Sequences
\newcommand{\seqfin}[3]{\ensuremath{#1_{#2}, \dotsc , #1_{#3}}}
\newcommand{\seqinf}[3]{\ensuremath{#1_{#2}, #1_{#3}, \dotsc}}

% New environments


% Header and footer
% Use the fancydr package
%\ifdraft{%
%\fancypagestyle{plain}{% redefine plain pagestyle for draft
%\renewcommand{\headrulewidth}{0pt}% no head rule in draft mode
%\fancyhf{}% clear headers and footers
%\fancyhead[C]{\ifdraft{DRAFT}{}}% print DRAFT across top if the draft option is set
%\fancyfoot[C]{\thepage}% usual page numbering
%

%\pagestyle{fancy}
%\renewcommand{\chaptermark}[1]{\markboth{\thechapter.\ #1}{}}% compact with no capitalisation
%\renewcommand{\sectionmark}[1]{\markright{\thesection.\ #1}}% compact with no capitalisation
%\ifdraft{\renewcommand{\headrulewidth}{0pt}}{}% no head rule in draft mode
%\setlength{\headheight}{15pt}
%\fancyhf{}% clear all header and footer fields
% one-sided printing, so no E,O distinction need be made
%\fancyhead[C]{\ifdraft{DRAFT}{}}% print DRAFT across top if the draft option is set
%\fancyhead[R]{\thepage}% page in top right
%\fancyhead[L]{\ifdraft{}{\leftmark}}% chapter number and title in top left

% Theorems
\theoremstyle{theorem}
\newtheorem{prop}{Proposition}[section]
\newtheorem{lemma}[prop]{Lemma}

\newtheorem{theorem}[prop]{Theorem}
%\newtheorem{hyp}[prop]{Hypothesis}
\newtheorem{cor}[prop]{Corollary}
\newtheorem{claim}[prop]{Claim}
%\newtheorem{notation}[prop]{Notation}
\theoremstyle{definition}

\newtheorem{defn}[prop]{Definition}
\newtheorem{assumption}[prop]{Assumption}
\newtheorem{eg}[prop]{Example}
\newtheorem{remark}[prop]{Remark}

% Numbering
\numberwithin{equation}{section}% standard style numbering, nothing special here

% Algorithms

% Use the algorithms package.

% \renewcommand{\algorithmicrequire}{\textbf{Input:}}
% \renewcommand{\algorithmicensure}{\textbf{Output:}}
% \renewcommand{\algorithmiccomment}[1]{/* #1 */}
% \RestyleAlgo{boxruled}
% \SetAlgoInsideSkip{medskip}
% \setlength{\algomargin}{2em}
% \setlength{\interspacetitleboxruled}{0.7em}
% \setlength{\interspacetitleruled}{0.7em}
% \LinesNumbered
% \SetAlCapNameSty{sc}
% \SetFuncSty{sc}

% \newenvironment{spacedalgorithm}{\begin{algorithm}\onehalfspacing}{\end{algorithm}}
% \newenvironment{onehalfverbatim}{\onehalfspacing\begin{verbatim}}{\end{verbatim}}

% \reversemarginpar
% \newcounter{nootje}
% \setcounter{nootje}{1}
% \renewcommand\check[1]{[*\thenootje]\marginnote{\tiny\begin{minipage}{40mm}\begin{flushright}\thenootje
% : #1\end{flushright}\end{minipage}}\addtocounter{nootje}{1}}

 
% Defining the custom listings environment
\definecolor{backcolour}{rgb}{0.95,0.95,0.92} % Background colour
\definecolor{codegray}{rgb}{0.5,0.5,0.5}
\definecolor{applegreen}{rgb}{0.55,0.71,0} % color for keywords in python 

\lstdefinestyle{gcstyle}{
    backgroundcolor=\color{backcolour},
    commentstyle=\color{codegray},
    keywordstyle=\color{applegreen},
    basicstyle=\ttfamily\footnotesize,
    breakatwhitespace=false,
    breaklines=true,
    keepspaces=true,
    captionpos=b,
    showspaces=false,
    showstringspaces=false,
    showtabs=false,
}

\lstset{style=gcstyle}

% Plotting trajectory tables 
\newlength{\figwidth}
\setlength{\figwidth}{0.40\textwidth}

\newenvironment{sidebyside}
{   
    \begin{tabular}{>{\centering\arraybackslash}p{2.5in}>{\centering\arraybackslash}p{2.5in}}   
}{%

\end{tabular}
}

% Parameter tables 
\newenvironment{paramtable}{
    \begin{tabular}{ccc}
        \hline 
        Parameter & Value & Description \\ \hline
}{%
    \end{tabular}
}

\newenvironment{summarytable}[1]{
    \begin{tabular}{cccccccc}
    \hline
    #1 & Mean & std & Min & Max & 25\% & 50\% & 75\% \\ 
    \hline
}{% 
    \end{tabular}

}


